\section{The \textsc{GAINS} Algorithm and Theoretical Analysis}
\label{sec:algorithm}

\textbf{\textsc{GAINS}} (\textbf{G}lobal \textbf{A}daptive \textbf{I}nference-time \textbf{N}oise \textbf{S}cheduling) operates at two granularities: offline coarse allocation (\cref{sec:offline}) sets per-timestep budgets from a small calibration run; online control (\cref{sec:online}) adapts them per prompt. The theory (\cref{sec:theory}) shows the per-step gain factorizes into a sensitivity factor times a sample-size factor, giving a closed-form water-filling allocation; an $\Omega(T)$ worst-case regret lower bound motivates online adaptation. Extensions to $\epsilon$-greedy/local-perturbation operators appear in \cref{sec:appendix-general}.

\subsection{Offline Coarse Time Scheduling}
\label{sec:offline}

Offline scheduling decides, before any test-time prompt arrives, how many local-search iterations to spend at each timestep. We consider the common setting where the sampler keeps a single latent $x_t$ and a single noise sample $\varepsilon_t$ at each timestep but performs multiple local search iterations before committing to $x_{t-1}$. The global scheduling problem then reduces to per-timestep decisions: at each $t$, given the search history and remaining budget, should we continue local search or stop and move on?

The difficulty is that, when deciding at $t$, we know neither (1) the additional improvement that further search at $t$ would bring nor (2) the improvement we forgo by spending the same NFE at later timesteps. Offline profiling addresses (1) and (2) jointly by estimating, for each timestep, how responsive the verifier score is to local search. For each model and dataset, we run a one-time profiling pass on a small calibration set: noise search executes with a uniform iteration count at every timestep, logging verifier scores after each iteration. From these logs we compute, for each $t$, the average per-iteration score gain and its variability across examples; aggregating reveals which segments of the trajectory respond most to local search. The pattern is strongly model-dependent: for Stable Diffusion, gains concentrate on early timesteps; for EDM they cluster in a contiguous middle segment (\cref{sec:appendix-experiments}). We partition timesteps into a \emph{high-value} region and a \emph{low-value} region, allocate a larger share of $B$ to the high-value region, and derive a coarse per-step budget $\{K_t\}$ (more iterations in high-value steps, fewer in low-value steps). The resulting schedule is tailored to the model and dataset but fixed across prompts; it captures \emph{where on the trajectory search matters most}.

\subsection{Online Fine-Grained Scheduling}
\label{sec:online}

Online scheduling refines, at test time, how the offline budget $\{K_t\}$ is spent on each individual prompt. At timestep $t$, local search runs up to $K_t$ iterations, but the controller can stop early or, within a small window, shift spending across timesteps. We distinguish two scores: $\hat s_t^{(j)}$, the raw verifier score of the candidate produced in iteration $j$, and $s_t^{(j)} := \max(s_t^{(j-1)}, \hat s_t^{(j)})$, the running best after iteration $j$ (with $s_t^{(0)} := \hat s_t^{(1)}$ from the first candidate). The controller monitors two statistics: the \emph{score gain} $g_t^{(j)} := s_t^{(j)} - s_t^{(j-1)} \ge 0$ (which is $0$ exactly when iteration $j$ failed to improve), and the \emph{score variance} across the candidates seen at step $t$,
\begin{equation}
  \mathrm{Var}_t^{(j)} := \mathrm{Var}_{i \le j}\bigl(\hat s_t^{(i)}\bigr).
\end{equation}

The two statistics target complementary questions. The score gain $g_t^{(j)}$ asks whether the current iteration is \emph{close to its local-search optimum}: when recent $g_t^{(j)}$ approach zero, the running best has converged and additional NFE buys little. The score variance $\mathrm{Var}_t^{(j)}$ asks whether this timestep is \emph{important}: by \cref{thm:loc-scale}(ii), within-step variance equals $\sigma_t^2$ to leading order, so $\mathrm{Var}_t^{(j)}$ estimates the per-step sensitivity driving every gain in \cref{sec:theory}. Stopping early is justified only when \emph{both} are small: a small gain alone can mark a high-sensitivity step waiting for a lucky draw, and small variance alone a low-sensitivity step whose iterations still nudge the running best upward; the conjunction prevents premature stopping in either regime (\cref{fig:stopping}). The \emph{revert-on-negative} construction (running max) makes $s_t^{(j)}$ monotone, so any iteration is at worst a no-op. To match realized compute exactly to $B$, over-/under-spending from adaptive stopping in the high-value region is offset against low-value steps. \cref{alg:gains} gives the procedure.

\begin{algorithm}[t]
\caption{\textsc{GAINS}: Global Adaptive Inference-time Noise Scheduling}
\label{alg:gains}
\begin{algorithmic}[1]
\STATE \textbf{Input:} timesteps $\{t_T,\dots,t_1\}$, total budget $B$, offline allocation $\{K_t\}$ ($\sum_t K_t = B$), slack $\delta$, gain window $W_g$, threshold coefs $(\beta_g,\beta_\sigma)$, black-box $\textsc{LocalIter}(t)$ that returns one raw candidate score $\hat s$ and within-step variance $\sigma$.
\STATE \textbf{Output:} realized per-step iterations $\{\widehat K_t\}$.
\STATE $R \leftarrow B$;\ \ $\mathcal{H}_g, \mathcal{H}_\sigma\leftarrow\emptyset$ (history).
\FOR{$t=T,T\!-\!1,\dots,1$}
    \STATE $K \leftarrow \min(K_t, R)$;\ $K_{\max}\leftarrow \min(K+\delta, R)$;\ $j_{\text{watch}}\leftarrow \max(2, K-\delta)$;\ $\mathcal{H}\leftarrow\emptyset$;\ $\widehat K_t \leftarrow 0$.
    \STATE $(\hat s^{(1)}, \sigma^{(1)}) \leftarrow \textsc{LocalIter}(t)$;\ $s^{(1)} \leftarrow \hat s^{(1)}$;\ $\widehat K_t \leftarrow 1$;\ $R \leftarrow R-1$.
    \FOR{$j = 2,\dots,K_{\max}$}
        \STATE $(\hat s^{(j)}, \sigma^{(j)}) \leftarrow \textsc{LocalIter}(t)$
        \STATE $s^{(j)} \leftarrow \max(s^{(j-1)}, \hat s^{(j)})$ \hfill {\small\COMMENT{running best (revert-on-negative)}}
        \STATE push $g^{(j)}\!\leftarrow\!s^{(j)}\!-\!s^{(j-1)} \ge 0$ into $\mathcal{H}$ (keep last $W_g$);\ $\widetilde g^{(j)}\!\leftarrow\!\mathrm{Mean}(\mathcal{H})$;\ $\widehat K_t\!\leftarrow\!\widehat K_t+1$;\ $R \leftarrow R-1$.
        \IF{$R=0$} \STATE \textbf{break} \ENDIF
        \IF{$j \ge j_{\text{watch}}$ \AND $|\mathcal{H}_g|>0$ \AND $|\mathcal{H}_\sigma|>0$ \AND $\widetilde g^{(j)} < \beta_g \mathrm{Mean}(\mathcal{H}_g)$ \AND $\sigma^{(j)} < \beta_\sigma \mathrm{Mean}(\mathcal{H}_\sigma)$}
            \STATE \textbf{break}
        \ENDIF
    \ENDFOR
    \STATE append $\mathrm{Mean}(\mathcal{H})$ to $\mathcal{H}_g$;\ append $\sigma^{(\widehat K_t)}$ to $\mathcal{H}_\sigma$.
\ENDFOR
\STATE \textbf{return} $\{\widehat K_t\}$.
\end{algorithmic}
\end{algorithm}

\subsection{Theoretical Foundations}
\label{sec:theory}

\emph{Scope.} The main-text results treat the local operator as \textbf{random search} ($K$ i.i.d.\ Gaussian candidates per timestep), giving the cleanest factorization $\sigma_t \cdot a_K$; the same form extends to $\epsilon$-greedy and local perturbation with operator-specific $\phi_K^{\mathcal L}$ replacing $a_K$ in \cref{sec:appendix-general}. We justify \textsc{GAINS} in three steps.
(1)~A Taylor expansion of the verifier score around the noiseless Euler point gives $S_t = \mu_t + \sigma_t Z_t + O(g_t^2 h)$ with $Z_t \sim \mathcal N(0,1)$, where $\sigma_t$ is the per-step \emph{sensitivity} that governs how much one extra noise candidate is worth, and the expected $K$-candidate gain is $\sigma_t a_K$ to leading order (\cref{thm:loc-scale}).
(2)~When $\sigma_t$ depends only on $t$, the optimal allocation has closed-form water-filling structure (\cref{thm:offline}).
(3)~When $\sigma_t = \sigma_t(c, x_t)$ varies per prompt, no adaptive policy avoids worst-case regret of order $T$ (\cref{lem:lower}; operator-agnostic). On non-adversarial prompt distributions, however, a positive Jensen gap (\cref{rmk:online-value}) leaves per-prompt online control valuable.\footnote{Notation: $g_t$ (diffusion coefficient, below) differs from $g_t^{(j)}$ (per-iteration score gain, \cref{sec:online}); $\sigma_t$ (sensitivity, below) differs from $\sigma$ (within-step empirical std, \cref{alg:gains}).}

\paragraph{SDE setup.}
We treat the DDPM update~\cref{eq:ddpm-reverse} as one Euler-Maruyama step~\citep{song2021scorebased} of a reverse-time SDE:
\begin{equation}
  X_{t-h} = \bar X_{t-h} + g_t\sqrt h\,\xi,\quad \bar X_{t-h} := X_t + b_t(X_t,c)\,h,\quad \xi\sim\mathcal N(0,I),
  \label{eq:em-step}
\end{equation}
with $h$ the integration step, $b_t$ the drift, $g_t$ the diffusion coefficient (both schedule-fixed), $\bar X_{t-h}$ the noiseless Euler point, and $\xi$ the per-step noise the local operator searches over. Let $D_\theta$ be the one-step denoiser~\citep{song2021scorebased}, $\Psi_t(x;c) := v(D_\theta(x,t{-}h,c);c)$ the verifier score on the one-step denoising, and $S_t := \Psi_t(X_{t-h};c)$. $\nabla\Psi_t$ denotes the gradient in the first argument at $\bar X_{t-h}$.

\begin{assumption}[Verifier smoothness]
\label{asm:smooth}
$\Psi_t(\cdot\,;c)\in C^2(\mathbb{R}^d)$, and there exists $L_t<\infty$ such that $\|\nabla^2\Psi_t(x;c)\|_{\mathrm{op}}\le L_t$ for all $x \in \mathbb{R}^d$.
\end{assumption}

\begin{theorem}[Location-scale structure of per-step scores]
\label{thm:loc-scale}
Under \cref{asm:smooth}, define $\mu_t := \Psi_t(\bar{X}_{t-h};c)$ and $\sigma_t := g_t\sqrt{h}\,\|\nabla\Psi_t(\bar{X}_{t-h};c)\|$. Conditional on $(X_t,c)$:
\begin{enumerate}
\item[\textup{(i)}] there exists $Z_t\sim\mathcal{N}(0,1)$ (the projection of $\xi$ onto the unit gradient direction $\nabla\Psi_t(\bar X_{t-h};c)/\|\nabla\Psi_t(\bar X_{t-h};c)\|$ when $\sigma_t>0$; arbitrary when $\sigma_t=0$, in which case the linear term vanishes) such that $S_t = \mu_t + \sigma_t Z_t + R_t$ with remainder $|R_t|\le\tfrac{1}{2}L_t g_t^2 h\|\xi\|^2$ a.s.\ ($Z_t$ and $R_t$ are not in general independent);
\item[\textup{(ii)}] as $g_t\sqrt{h}\to 0$, $\mathrm{Var}(S_t\mid X_t,c) = \sigma_t^2 + O_{d, L_t, \|\nabla\Psi_t\|}(g_t^3 h^{3/2})$, where the implicit constant depends on dimension $d$, the local Hessian bound $L_t$, and the gradient norm at $\bar X_{t-h}$;
\item[\textup{(iii)}] for $K$ i.i.d.\ noise candidates $\xi^{(k)}\stackrel{\text{iid}}\sim\mathcal N(0,I)$ at step $t$ (random search), let $S_t^{(k)}:=\Psi_t(\bar X_{t-h}+g_t\sqrt h\,\xi^{(k)};c)$. Then
\begin{equation}
G_t(K) := \mathbb{E}\Bigl[\max_{1\le k\le K} S_t^{(k)}\,\Big|\,X_t,c\Bigr] - \mu_t = \sigma_t\,a_K + O_{K,d,L_t}(g_t^2 h),
\label{eq:gain-sde}
\end{equation}
where $a_K := \mathbb E[\max(Z_1,\ldots,Z_K)]$, $Z_k\overset{\text{iid}}\sim\mathcal N(0,1)$, with $a_K\sim\sqrt{2\log K}$.
\end{enumerate}
\end{theorem}

The proof (\cref{sec:proof-loc-scale}) is a second-order Taylor expansion of $\Psi_t$ around $\bar{X}_{t-h}$ (\cref{prop:taylor}). The product $\sigma_t = g_t\sqrt{h}\,\|\nabla\Psi_t\|$ of diffusion coefficient, step size, and verifier gradient norm is the leading-order coefficient governing how much extra search at $t$ is worth.

\paragraph{Optimal allocation: water-filling.}
Treating $\{\sigma_t\}$ as a fixed input and dropping the $O(g_t^2 h)$ remainder in \cref{eq:gain-sde}, summing over $t$ gives the surrogate
\begin{equation}
  \max_{\{K_t\}_{t=1}^T}\;\sum_{t=1}^{T}\sigma_t\,a_{K_t}
  \quad\text{s.t.}\;\sum_{t=1}^{T}K_t=B,\;K_t\in\mathbb{Z}_{\ge 1}.
  \label{eq:alloc}
\end{equation}
Write $\Delta a_K := a_{K+1}-a_K$ for the discrete marginal of $a_K$, and call $t$ \emph{active} at $\{K_t^*\}$ if $K_t^* > 1$.

\begin{theorem}[Offline optimality under independence]
\label{thm:offline}
Assume $\sigma_t$ depends only on $t$ (not on $c$ or $x_t$) and the remainder $O_{K,d,L_t}(g_t^2 h)$ in \cref{thm:loc-scale}\textup{(iii)} is uniform in $(c,x_t)$. Then the optimizer $\{K_t^*\}$ of \cref{eq:alloc} satisfies:
\begin{enumerate}
\item[\textup{(i)}] \emph{Prompt-independence.} The same $\{K_t^*\}$ optimizes \cref{eq:alloc} for every $(c, x_t)$.

\item[\textup{(ii)}] \emph{Discrete KKT sandwich.} There exists $\lambda^* \ge 0$ (the dual variable for $\sum_t K_t = B$) such that, for every active $t$,
\begin{equation}
  \sigma_t\,\Delta a_{K_t^* - 1} \;\ge\; \lambda^* \;\ge\; \sigma_t\,\Delta a_{K_t^*}.
  \label{eq:kkt-sandwich}
\end{equation}

\item[\textup{(ii')}] \emph{Weak monotonicity.}
\begin{equation}
  \sigma_{t_1} > \sigma_{t_2} \;\Longrightarrow\; K_{t_1}^* \ge K_{t_2}^*.
\end{equation}
At the integer level the inequality is strict iff $\sigma_{t_1}/\sigma_{t_2} > \Delta a_{K-1}/\Delta a_K$ at the binding $K$; the continuous relaxation $\tilde a \in C^1$ of $a_K$ gives strict monotonicity unconditionally.

\item[\textup{(iii)}] \emph{Dominance over uniform.} If $T \mid B$ so $K_t \equiv B/T$ is feasible,
\begin{equation}
  \sum_{t=1}^T \sigma_t\,a_{K_t^*} \;\ge\; \sum_{t=1}^T \sigma_t\,a_{B/T},
  \label{eq:dominance}
\end{equation}
with strict inequality in the continuous relaxation whenever $B > T$ and $\sigma_{t_1}\ne\sigma_{t_2}$ for some pair (at the integer level, conditional on the gap condition in (ii')). Defining the across-timestep dispersion
\begin{equation}
  \mathrm{Var}(\sigma_t) := \frac{1}{T}\sum_{t=1}^T (\sigma_t - \bar\sigma)^2, \qquad \bar\sigma := \frac{1}{T}\sum_t \sigma_t,
  \label{eq:dispersion}
\end{equation}
the first-order improvement along any budget-neutral pair perturbation $K_{t_1}\to B/T+\epsilon,\,K_{t_2}\to B/T-\epsilon$ is $\epsilon(\sigma_{t_1}-\sigma_{t_2})\,\tilde a'(B/T)+O(\epsilon^2)$ (we use $\mathrm{Var}(\sigma_t)$ as a summary dispersion measure, without claiming global monotonicity). The case $B=T$ forces $K_t\equiv 1$.
\end{enumerate}
\end{theorem}

The proof (\cref{sec:proof-offline}) is by Lagrangian first-order conditions and concavity of $a_K$. The same water-filling structure transfers to any local operator $\mathcal L$ satisfying the regularity conditions of \cref{asm:local-search}, with $a_K$ replaced by an operator-specific gain $\phi_K^{\mathcal L}$ (\cref{sec:appendix-general}).

\paragraph{Regret bound for adaptive scheduling.}
When $\sigma_t = \sigma_t(c, x_t)$ depends on the prompt, $\{K_t^*\}$ is no longer per-instance optimal. Embed the problem into an abstract sequential-allocation game: at each round $t$ the learner sees the current reward distribution $\mathcal{D}_t$ (in the diffusion problem, the law of $\sigma_t a_{K_t}$), picks $K_t \in \{0,1\}$ subject to $\sum_t K_t \le B$, and gets reward $\sim\mathcal{D}_t$ if $K_t = 1$ (else $0$). Regret is defined against the offline benchmark that knows $\mathcal{D}_1,\ldots,\mathcal{D}_T$ in advance.

\begin{lemma}[Regret lower bound for the sequential-allocation game]
\label{lem:lower}
For every (possibly randomized) adaptive policy there exists an instance with $T = 2B$ such that
\begin{equation}
  \mathbb{E}[R_T] \;\ge\; \tfrac{B}{4} \;=\; \Omega(T).
  \label{eq:lower}
\end{equation}
\end{lemma}

A diffusion policy with prompt-dependent $\sigma_t$ simulates this game by treating each step as a 0/1 commit, so $\Omega(T)$ regret transfers. The construction (\cref{sec:proof-lower}) sets $\mathcal{D}_t = \delta_{1/2}$ for $t \le B$ and either $\delta_0$ or $\delta_1$ for $t > B$, making the two instances statistically identical through round $B$. The bound applies to adversarial $\boldsymbol\sigma$; \cref{rmk:online-value} below governs the stochastic-$\boldsymbol\sigma$ regime targeted by \textsc{GAINS}.

\begin{remark}[Online controller as Jensen-gap recovery]
\label{rmk:online-value}
For prompt-dependent $\boldsymbol\sigma = \boldsymbol\sigma(c, x_t)$, define the static oracle value
\begin{equation}
  V^*(\boldsymbol\sigma) := \max_{\substack{\sum_t K_t = B\\ K_t \ge 1}}\;\sum_t \sigma_t\,a_{K_t}.
  \label{eq:static-oracle}
\end{equation}
$V^*$ is convex in $\boldsymbol\sigma$ (pointwise maximum of linear functions), so by Jensen's inequality
\begin{equation}
  J(\boldsymbol\sigma) \;:=\; \mathbb{E}_{\boldsymbol\sigma}[V^*(\boldsymbol\sigma)] - V^*(\bar{\boldsymbol\sigma}) \;\ge\; 0,
  \label{eq:jensen-gap}
\end{equation}
where $\bar{\boldsymbol\sigma} := \mathbb{E}\,\boldsymbol\sigma$. We call $J(\boldsymbol\sigma)$ the \emph{Jensen gap}: the expected gain of an oracle that allocates per prompt over one tuned to the population-mean profile. \cref{lem:lower} prevents closing $J$ in the adversarial worst case, but structured $\boldsymbol\sigma$ admits partial recovery. \textsc{GAINS} stops at step $t$ only when both the recent gain $\widetilde g^{(j)}$ and the within-step variance $\sigma^{(j)}$ are small (verifier locally insensitive \emph{and} search no longer improving); either signal alone keeps search active. \cref{fig:stopping} illustrates the rule. The convexity argument carries over to general operators (\cref{sec:appendix-general}).
\end{remark}

\begin{figure}[t]
\centering
\includegraphics[width=0.45\linewidth]{figures/fig_stopping_decision.pdf}
\caption{Online stopping decision in the $(\sigma_t, \Delta a_K)$ plane. Horizontal threshold: small recent gain; vertical threshold: small empirical variance (a sensitivity proxy). Stopping triggers only when \emph{both} fire (shaded). The hyperbola $\sigma_t\cdot\Delta a_K=\lambda^*$ is the true optimality boundary; the rectangular proxy is conservative.}
\label{fig:stopping}
\end{figure}
